\section{行省接收下的不同地方}

1267年营建大都，1277年开始江南行省、路府和财赋接收，1281年对日远征失利，1344年河患加重北方压力，1351年河工与红巾军起事并见。它们不宜被串成“中央强弱”的单线故事：大都的物资、江南的征收、海域的造船与北方河工，分别依赖不同的工匠、站赤、仓储、河道和地方合作。

《元典章》《通制条格》与行省文书能说明法令和机构，不能保证实际执行。纸钞数额、户籍分类、站户或匠户的规定，都须与碑刻、契约、工程遗存和区域材料相互检验；“色目”“汉人”等史料分类更不能代替具体的职业、地域与法律身份。元末的灾害和战争使这种不平衡加剧，却不能由一场河工或一次起事解释全部地方社会。
